\begin{abstract}
IO memory protection prevents malicious DMA-capable IO devices from executing errant data transfers 
    into host memory. 
Modern servers enable memory protection by requiring IO devices 
    to operate on virtual IO addresses, 
    and using IOMMU % (along with specialized caches) 
    to perform fast virtual-to-physical address translation upon each DMA. 
Existing IO memory protection mechanisms offer powerful, so-called strict, safety guarantees; 
however, they require at least one VM exit per DMA in virtualized environments 
    to invalidate translation entries in IOMMU caches. 
These VM exits in strict protection---even with state-of-the-art 
    device passthrough mechanisms, memory management, 
    and hardware-assisted translation---result in 83--95$\%$ 
    degradation in application-level throughput! 
Real-world deployments today have to choose either high performance 
    or strict IO memory protection, but not both. 

We present {\name}, a simple modification to existing IO memory protection mechanisms 
    that makes strict IO memory protection feasible in virtualized environments. 
The key insight in {\name} is that, in virtualized environments, 
    it is possible to reduce the overheads of strict IO memory protection 
    to a bare minimum by \textit{increasing} the cost of address translation. 
Intuitively, this is because each VM exit has much higher cost 
    than the worst-case translation cost. 
{\name} design incorporates two ideas---\textit{one-for-many invalidations} and \textit{deferred free page}---that, 
    when put together, result in higher address translation cost but significantly 
    reduce the number of VM exits (and hence, IO memory protection overheads) 
    while maintaining strict IO memory protection. 
We implement {\name} on Linux, and evaluate it across multiple real-world applications. 
Our evaluation suggests that, at the cost of one dedicated core, 
    {\name} consistently achieves performance within 1--8$\%$ of 
    the unprotected system---74--147$\times$ higher than state-of-the-art IO memory protection mechanisms. 

% We consider the problem of IO memory protection in virtualized environments.
% I/O memory protection in virtual machines (VMs)
%   defends guest systems against errant DMAs
%   through virtualized IOMMUs (vIOMMUs).
% However, ensuring strong isolation for data integrity and confidentiality 
%     requires invalidating I/O address translations 
%     from hardware caches (e.g., IOTLBs) after every DMA completion.
% Since such invalidations must be 
%   done by the host, they cause frequent VM exits
%   and incur excessive performance overhead
%   (up to 95\% throughput drops). 
% Consequently, existing systems often relax protection guarantees, 
%     yet still suffer non-trivial performance costs.

% This paper shows that the cost of I/O memory protection
%   in virtualized environments is not inherent.
% We present \name, a new design that delivers strict vIOMMU protection 
%   at near-zero overhead compared to unprotected DMA. 
% The insight is that frequent VM exits for hardware invalidation are not necessary but
%   artifacts of 
%   globally serializing invalidation across threads
%   to enforce safety.
% \name{} eliminates all unnecessary serializations that are not essential
%   to the safety invariant of strict protection which only requires 
%   strict ordering of invalidations {\it per thread}.
% It promotes highly concurrent invalidation processing 
%   across threads,
%   avoiding unnecessary VM exits and scaling with modern multi-core systems,
% without compromising isolation.
% % \name{} is realized by two complementary techniques:
% % \emph{safely batched invalidation} which decouples the preparation 
% %  of invalidation requests from their submission, which offers 
% %  the opportunities to effectively coalescing concurrent requests,
% % and \emph{deferred free pages} which enforces safety ordering
% %  without blocking the CPU during invalidation.
% % We implement \brand in Linux kernel 6.12.9 with fewer than 1.5k lines of change.
% On modern hardware with nested IOMMU and 400 Gbps networking,
%   \name\ achieves up to 7$\times$ throughput improvement on Rx traffic
%   and up to 42$\times$ on Tx traffic over vanilla Linux/KVM
%   with strict protection; \name's performance is close to that without 
%   using vIOMMU.

\end{abstract}